\section{Formal verification of the certificate theorem}
\label{sec:formal-verification}

A Lean development verifies Theorem~\ref{thm:certificate} and its HHC
specialization, Corollary~\ref{cor:hhc-certificate}, for finite families
of real symmetric matrices, real vectors, and real constants. Its strict
feasible set, ordinary convex hull, and nontrivial certificate agree with
those defined here. The main declaration is
\texttt{proper\_hull\_iff\_certificate} in the namespace
\texttt{QuadraticAggregation.System}; the HHC specialization is
\texttt{proper\_hull\_iff\_certificate\_of\_hhc}.
The assumptions are nonemptiness and AHC, or nonemptiness and HHC,
respectively. No positive definite combination is assumed.
The certificate-to-proper-hull implication is proved separately without
either assumption.

The formal definitions express HHC through kernels of nonzero linear
functionals, which are precisely the linear hyperplanes in this
finite-dimensional setting. AHC is stated using arbitrarily large good
levels, as in Definition~\ref{def:ahc}; its equivalence with the increasing
sequence formulation is also proved. Matrix positive semidefiniteness is
linked to nonnegativity of the quadratic form, using the symmetry of the
original matrices and their aggregations.

The proof follows the constant-elimination and coefficient-normalization
argument in Section~\ref{sec:certificate}. In particular, it constructs
separation weights from the convex hyperplane images and obtains the
nonzero limit in the actual finitely generated coefficient cone. These
are proved steps, not additional hypotheses of the headline theorem.
The implementation uses the entrywise maximum matrix norm and the
product maximum norm for coefficient pairs. These are equivalent to the
Frobenius and Euclidean norms in finite dimension; the compactness proof
works in either choice.

The verified supporting results include
Lemmas~\ref{lem:negative-direction}, \ref{lem:sweeping},
\ref{lem:hyperplane-certificate}, and~\ref{lem:finite-cone}, and the
uniform cone-separation lemma, Lemma~\ref{lem:uniform-cones}.
The last lemma is checked independently, although the main formal proof
does not use it. The spectral distance calculation and quantitative
alternative in Section~\ref{sec:quantitative-obstruction} are not
formalized as separate results.

This core package consists of eleven modules in
\path{supplement/lean/Formal/QuadraticAggregation/}, accompanied by a declaration
coverage map and reproducible checks. With the pinned Lean 4.33.1
toolchain, targeted compilation without warnings, a transitive axiom
audit of all 178 declarations owned by those modules, and kernel replay
of each module passed. The checks were rerun during preparation of this
manuscript, with source fingerprints unchanged throughout the run.
The only permitted axioms were \texttt{propext},
\texttt{Classical.choice}, and \texttt{Quot.sound}.
Kernel replay uses the pinned Lean kernel and imported dependencies;
it does not constitute verification by an independently implemented
proof assistant or a replay of all dependencies from source.

This formal scope concerns existence of a nontrivial globally convex
aggregation. It does not establish that such aggregations describe the
entire hull. The closed-system and relaxation consequences, SDP
characterizations, examples, applications, further aggregation results,
and literature-priority claims are outside this core package. Separate
packages verify the specified later results described in the accompanying
formal accounts.
Formal verification of the stated theorem does not certify the rest
of the manuscript.
